\section{Sprint candidate 3: Inheritance-tax reforms and outcomes}
\label{sec:tax}

\subsection{Five-day question}
What changes in firm births and house prices can be detected around six US state estate/inheritance-tax repeals, and how sensitive are estimates to reform timing, comparison groups and statistical power?

\subsection{Data package by December 2026}
Assemble a public panel of 50 states plus DC, targeting 2003--2024 (up to 1,122 state-years, with source-specific missingness), with reforms in Ohio, North Carolina, Indiana, Tennessee, New Jersey and Delaware (2013--2018).

\begin{itemize}[nosep]
\item \textbf{Primary outcomes:} Census Business Dynamics Statistics (firm/establishment births) and FHFA house-price indices.
\item \textbf{Controls/secondary outcomes:} BEA state income, GDP and consumption; ACS tenure and demographics; New York Fed debt balances.
\item \textbf{Reform register:} announcement, passage and effective dates, tax type, rates/exemptions, affected estates, concurrent changes and source passages. Keep New Jersey's 2017 exemption rise distinct from its 2018 estate-tax repeal.
\item \textbf{Reform intensity:} state death/gift tax collections (Census T50) before repeal, with documented limitations.
\end{itemize}

Prepare tidy data, a variable dictionary, verified legislation and runnable baseline notebooks. This package uses aggregate public data and has no beneficiary-identity access dependency.

\subsection{Five-day deliverables}
Staggered event studies/difference-in-differences with confidence intervals, pre-trends, anticipation and placebo checks, and a power assessment. The first outcomes are business formation and house prices; consumption and debt are extensions if ready. A separate optional task tests AI extraction of policy parameters against the checked register.

\subsection{To do by 31 December 2026}
\begin{itemize}[nosep]
\item Download, clean and merge BDS, FHFA and BEA core series; document units and coverage.
\item Double-check repeal dates and legislative steps against primary sources.
\item Build treatment, comparison-state and concurrent-policy variables.
\item Run baseline estimates and a power check; freeze a reproducible public data package.
\item Add optional NY Fed/ACS series and AI document tests only after the core works.
\item Make sure data sprint participants have access to the approved datasets, documentation and required computing environment before the sprint.
\item Prepare documents for data sprint participants explaining the tasks, suggested teams, division of work, coordination points and possible expansions or stretch goals.
\end{itemize}

\subsection{Participant preparation and suggested teams}
\begin{description}[style=nextline]
\item[Access preflight] Test the public US panel, reform register, legislative files and reproducible baseline notebooks in the actual participant environment; record data redistribution terms.
\item[Participant brief] Provide a dated reform chronology, variable dictionary, comparison-state candidates and a short explanation of confounding, staggered treatment and statistical power.
\item[Suggested teams] Team A validates tax events and concurrent policies; Team B harmonises outcomes and estimates event studies. Integrate register changes into sensitivity tests.
\item[Stretch goals] Synthetic controls, reform-size heterogeneity, international comparisons or AI extraction of legislative details.
\end{description}

\subsection{Longer-term extension}
Study additional outcomes, tax-reform magnitude and heterogeneous responses; add synthetic-control designs and country panels for Norway (2014), Finland (2013/2015/2017) and South Africa (2018), screened for confounding reforms. Sweden, Japan, Italy, the Philippines and Colombia are possible further policy-document cases. Administrative beneficiary-level effects require a separate research project.
